\chapter*{ประมวลรายวิชา}
\addcontentsline{toc}{chapter}{ประมวลรายวิชา}

\noindent\textbf{รหัสและชื่อรายวิชา:} เกษตรอัจฉริยะและปัญญาประดิษฐ์ในสวนไม้ผล (Smart Agriculture and AIoT for Orchard Systems)\\[0.2cm]
\noindent\textbf{จำนวนหน่วยกิต:} 3 (2-2-5) หน่วยกิต (ทฤษฎี 2 คาบ ปฏิบัติการ 2 คาบ ศึกษาค้นคว้าด้วยตนเอง 5 คาบต่อสัปดาห์)\\[0.2cm]
\noindent\textbf{อาจารย์ผู้รับผิดชอบรายวิชา:} ผู้ช่วยศาสตราจารย์ ดร.ชีวะ ทัศนา\\[0.2cm]
\noindent\textbf{สาขาวิชาและคณะ:} สาขาวิชาฟิสิกส์ คณะวิทยาศาสตร์และเทคโนโลยี มหาวิทยาลัยราชภัฏรำไพพรรณี

\section*{คำอธิบายรายวิชา}
หลักการและสถาปัตยกรรมของระบบเกษตรอัจฉริยะ การประยุกต์ใช้อินเทอร์เน็ตของสรรพสิ่งทางการเกษตร (AIoT) เซนเซอร์ตรวจวัดสภาพแวดล้อมทางกายภาพ ชีววิทยา และเคมีในดินและบรรยากาศ สถาปัตยกรรมเครือข่ายสื่อสารไร้สายระยะไกลกำลังต่ำ (LoRaWAN) ฟิสิกส์ทุเรียนและการสะสมชีวมวล ระบบการชลประทานอัจฉริยะแบบแม่นยำสูง การประยุกต์ใช้ปัญญาประดิษฐ์และการประมวลผลภาพถ่ายดิจิทัลสมาร์ทโฟนเพื่อการวิเคราะห์คุณภาพดิน (SoilAI) การบูรณาการแดชบอร์ดข้อมูลเพื่อการตัดสินใจ และกรณีศึกษาการจัดการแปลงไม้ผลเศรษฐกิจในจังหวัดจันทบุรี

\section*{วัตถุประสงค์การเรียนรู้ (Course Learning Outcomes; CLOs)}
เมื่อนักศึกษาได้ศึกษาตามหลักสูตรรายวิชานี้แล้ว จะต้องมีความรู้ความสามารถดังนี้:
\begin{enumerate}
    \item อธิบายหลักการทำงานของเซนเซอร์เกษตร เครือข่ายไร้สาย LoRaWAN และสถาปัตยกรรม AIoT ได้อย่างถูกต้อง
    \item วิเคราะห์พารามิเตอร์ทางกายภาพ ความชื้นในดิน อุณหภูมิ และแสงอาทิตย์ที่มีผลต่อพัฒนาการของผลทุเรียนได้
    \item ติดตั้ง กำหนดค่า และสอบเทียบอุปกรณ์เซนเซอร์วัดค่าในดินและบรรยากาศตามมาตรฐานวิชาการได้
    \item ใช้งานระบบปัญญาประดิษฐ์วิเคราะห์คุณภาพดินจากภาพถ่ายสมาร์ทโฟน และแปลผลการวินิจฉัยเพื่อการจัดการธาตุอาหารพืชได้อย่างมีประสิทธิภาพ
    \item ออกแบบผังระบบชลประทานอัตโนมัติและแดชบอร์ดแสดงผลข้อมูลสำหรับการจัดการสวนทุเรียนอัจฉริยะได้
\end{enumerate}

\section*{แผนการประเมินผลการเรียนรู้}
\begin{table}[H]
\centering\small
\begin{tabularx}{\textwidth}{p{0.55\textwidth} p{0.25\textwidth} c}
\toprule
\textbf{กิจกรรมการประเมินผล} & \textbf{สัปดาห์ที่ประเมิน} & \textbf{สัดส่วนคะแนน} \\
\midrule
1. การมีส่วนร่วมในชั้นเรียนและการตรงต่อเวลา & ตลอดภาคการศึกษา & 10\% \\
2. รายงานปฏิบัติการทดลองและการสอบเทียบเซนเซอร์ & สัปดาห์ที่ 3, 5, 7, 10 & 30\% \\
3. โครงงานออกแบบระบบ AIoT และแปลงทดสอบสวนทุเรียน & สัปดาห์ที่ 13--14 & 20\% \\
4. การสอบวัดผลกลางภาค & สัปดาห์ที่ 8 & 20\% \\
5. การสอบวัดผลปลายภาค & สัปดาห์ที่ 16 & 20\% \\
\midrule
\textbf{รวมทั้งสิ้น} & & \textbf{100\%} \\
\bottomrule
\end{tabularx}
\end{table}
